\section{Explicit centering and the height carrying half the total}
\label{sec:median}

We now specialize to the positive noncentral zero ordinates of $\zeta$
or of a fixed primitive real Dirichlet character $\chi$. For $\zeta$ put
$q=1$; otherwise $q$ is the conductor. Keep a fixed $\theta>0$.
Every occurrence of a zero is counted, so a zero of multiplicity $m$
contributes $m$ independent shape-$\theta$ variables, equivalently one
shape-$m\theta$ variable. A possible central zero is excluded from the
series and removed from the completed-function product.

The positive-ordinate counting formula is
\begin{equation}\label{med:eq:count}
 N(T)=\frac{T}{2\pi}\left(\log\frac{qT}{2\pi}-1\right)
       +O_\chi(\log(T+2)).
\end{equation}
For real characters this is half the usual both-sign count, up to the
finite central multiplicity; see~\cite{BennettEtAl}. For $\zeta$ it is
the Riemann--von Mangoldt formula. All constants below may depend on the
fixed function and on $\theta$.

Let $u$ be the exact saddle from \eqref{mc:eq:saddle} with $p=2$, and set
\[
 \ell=\log\frac{qu}{2\pi},\qquad
 J(c)=\int_0^c\frac{\log r}{1+r^2}\,dr,\qquad
 M_u(c)=\frac{\theta}{2\pi u}
                 \bigl(\ell\arctan c+J(c)\bigr).
\]
The first correction $J$ matters: $J(1)=-G$, where
$G=\sum_{n\ge0}(-1)^n/(2n+1)^2$ is Catalan's constant.

\begin{lemma}[Centering accurate on the fluctuation scale]
\label{med:lem:center}
Uniformly for $c$ in any compact subset of $(0,\infty)$,
\begin{align}
 m_x(c)&=M_u(c)+O_\chi\left(\frac{\log u}{u^2}\right)
                  =M_u(c)+o(\sigma_x),\label{med:eq:mean}\\
 x&=\frac{\theta\ell}{4u}
                  +O_\chi\left(\frac{\log u}{u^2}\right),\label{med:eq:total}\\
 \sigma_x^2&=\frac{\theta(\ell-1)}{8u^3}
                  +O_\chi\left(\frac{\log u}{u^4}\right).
                  \label{med:eq:variance}
\end{align}
\end{lemma}

\begin{proof}
Write the error in \eqref{med:eq:count} as $R(T)$, extending the formula
to $T\ge0$ by increasing its constant; then $R(0)=0$. Stieltjes
integration by parts with $f(T)=(T^2+u^2)^{-1}$ gives
\[
 m_x(c)=\frac{\theta}{2\pi u}
       (\ell\arctan c+J(c))
       +\theta f(cu)R(cu)-\theta\int_0^{cu}f'(T)R(T)\,dT.
\]
Both error terms are $O_\chi(\log u/u^2)$, uniformly on the stated
compact sets. The same estimate holds at $c=\infty$; the boundary term
there vanishes and the integral is absolutely convergent. Substitution
$r\mapsto1/r$ gives $J(\infty)=0$, proving \eqref{med:eq:total}.

Apply the same argument to $(T^2+u^2)^{-2}$. The main integrals are
\[
 \int_0^\infty\frac{dr}{(1+r^2)^2}=\frac\pi4,
 \qquad
 \int_0^\infty\frac{\log r}{(1+r^2)^2}\,dr=-\frac\pi4.
\]
For the second identity, use
$(1+r^2)^{-2}=\tfrac12(r/(1+r^2))'+\tfrac12(1+r^2)^{-1}$,
integrate by parts, and use $J(\infty)=0$. The error is
$O_\chi(\log u/u^4)$, giving \eqref{med:eq:variance}.
Finally
$(\log u/u^2)/\sigma_x=O(\sqrt{\log u/u})\to0$.
\end{proof}

\begin{remark}[A leading-order mean is insufficient]
The uniform law of large numbers permits the centering $xA(c)$ at scale
$x$, but the bridge theorem generally does not permit it at scale
$\sigma_x$. For instance
\[
 m_x(1)-\frac x2=-\frac{\theta G}{2\pi u}
                    +O_\chi\left(\frac{\log u}{u^2}\right).
\]
The magnitude of this deterministic difference divided by $\sigma_x$
grows on the order of $\sqrt{u/\log u}$. Exact or sufficiently accurate
centering is therefore a mathematical requirement.
\end{remark}

Define the random height
\begin{equation}\label{med:eq:height}
 H_x=\min\left\{\gamma_j:
        \sum_{\gamma_i\le\gamma_j}\frac{G_i}{\gamma_i^2}
                 \ge\frac X2\right\}.
\end{equation}
It exists almost surely because the partial sums increase to $X>0$.
For $\ell>0$, let $c_\ell>1$ be the unique solution of
\begin{equation}\label{med:eq:centerroot}
 \ell\arctan c_\ell+J(c_\ell)=\frac{\pi\ell}{4}.
\end{equation}
The equation is a scalar monotone root problem: its left side minus
the right side is $-G$ at $1$, tends to $\pi\ell/4$ at infinity, and has
positive derivative $(\ell+\log c)/(1+c^2)$ for $c\ge1$.

\begin{theorem}[Conditional half-contribution height]
\label{med:thm:height}
Under the conditional law $X\le x$, as $x\downarrow0$,
\begin{equation}\label{med:eq:CLT}
 \boxed{\quad
 \frac{x}{\pi\sigma_x}\left(\frac{H_x}{u}-c_\ell\right)
 \ \Longrightarrow\
 N\left(0,\frac14-\frac1{\pi^2}\right).
 \quad}
\end{equation}
The center and height scale satisfy
\begin{align}
 c_\ell&=1+\frac{2G}{\ell}+\frac{2G^2}{\ell^2}
                            +O(\ell^{-3}),\label{med:eq:expansion}\\
 \frac{u\pi\sigma_x}{x}
       &\sim\pi\sqrt{\frac{2u}{\theta\ell}}.
       \label{med:eq:heightscale}
\end{align}
In \eqref{med:eq:CLT} the implicit center \eqref{med:eq:centerroot}
must be used, not a fixed truncation of \eqref{med:eq:expansion}.
\end{theorem}

\begin{proof}
We first justify a moving-cutoff version of the one-band assertion in
Theorem~\ref{mc:thm:joint}. If $c_u\to1$, split the tilted coordinates
at $c_uu$. The variance fractions of the two centered blocks tend to
\[
 v_1=V(1)=\frac12+\frac1\pi,\qquad 1-v_1.
\]
This follows from the spectral sum lemma and monotone bracketing of
$c_u$ between fixed cutoffs. The local limit proof remains valid:
$[u/4,u/2]$ lies in the lower block and $[2u,3u]$ in the upper block
for all sufficiently large $u$. Both intervals contain order
$u\log u$ ordinates and supply the same integrable Fourier envelope.
Thus their normalized densities converge uniformly to
$\varphi_{v_1}$ and $\varphi_{1-v_1}$.

Applying the exact two-band density change of variables from
\eqref{mc:eq:jointdensity} proves
\begin{equation}\label{med:eq:moving}
 \left(\frac{X_x(c_u)-m_x(c_u)}{\sigma_x},
              u^2(x-X)\right)
 \Longrightarrow (Z,E),
\end{equation}
where $Z\sim N(0,v_1(1-v_1))$ and $E\sim\operatorname{Exp}(1)$ are
independent. This step is only a two-dimensional density argument;
it requires no functional central limit theorem.

Put $d_u=\pi\sigma_x/x\to0$ and fix $z\in\mathbb R$. Take
$c_u=c_\ell+zd_u$, which tends to one. Lemma~\ref{med:lem:center} and
\eqref{med:eq:centerroot} give $m_x(c_\ell)-x/2=o(\sigma_x)$. Also
\[
 M'_u(c)=\frac{\theta}{2\pi u}
                  \frac{\ell+\log c}{1+c^2}
          \sim\frac{x}{\pi}\qquad(c\to1),
\]
uniformly in a shrinking neighborhood of one. Hence
\[
 \frac{m_x(c_u)-x/2}{\sigma_x}\longrightarrow z.
\]
Under the conditional law,
$(x-X)/(2\sigma_x)=u^2(x-X)/(2u^2\sigma_x)\to0$ in probability,
by \eqref{med:eq:moving}. Consequently
\[
 \frac{X_x(c_u)-X/2}{\sigma_x}\Longrightarrow Z+z.
\]
The exact monotone equivalence
$H_x\le uc_u\Longleftrightarrow X_x(c_u)\ge X/2$ now shows that the
distribution function in \eqref{med:eq:CLT} tends at $z$ to
$\Pr(Z\ge-z)$, the stated normal distribution function. Repeated
ordinates cause no problem for this equivalence.

Since $v_1(1-v_1)=1/4-1/\pi^2$, the variance is as claimed.
Equations~\eqref{med:eq:total}--\eqref{med:eq:variance} yield
\eqref{med:eq:heightscale}. Finally, Taylor expansion at $c=1$, using
$J(1)=-G$ and $J'(1)=0$, gives \eqref{med:eq:expansion}.
\end{proof}

The second term omitted by the first-order center
$1+2G/\ell$ is asymptotic to $2G^2/\ell^2$. After the normalization
in \eqref{med:eq:CLT}, this error grows like
$\sqrt u/\ell^{3/2}$. Thus even the first correction, although larger
than the random fluctuations and informative about their location,
cannot replace the exact root in the central limit theorem.

